\section{Metodologia}\label{sec:metodologia}

Para compreender a evolução das Ligas de Alta Entropia em regimes hipersônicos, este estudo desenvolveu um levantamento bibliométrico na base de dados \emph{Web of Science Core Collection} abrangendo o período de 2004 a 2026. O objetivo central consistiu em estruturar um protocolo transparente e determinístico, assegurando que pesquisadores independentes possam replicar integralmente cada etapa da triagem e recuperar os mesmos registros documentais.

\subsection{Estratégia de Busca e Combinação Booleana}\label{sec:estrategia_busca}

A formulação da expressão de busca assemelha-se ao ajuste focal de uma lente óptica, onde a amplitude excessiva capta ruídos irrelevantes e a restrição severa descarta dados essenciais. Enquanto buscas genéricas por propriedades mecânicas em alta temperatura retornam milhares de publicações desvinculadas da severidade aeroespacial, a conjunção direta e exclusiva das palavras ``HEAs'' e ``hipersônico'' recupera apenas 38 estudos, omitindo marcos fundamentais da termodinâmica de ligas refratárias.

Para superar essa limitação, desenhou-se uma estratégia combinatória automatizada em linguagem Python baseada no cruzamento de duas categorias temáticas completas. A primeira categoria congrega 42 variações terminológicas dedicadas a ligas de alta entropia, composicionalmente complexas e ligas refratárias concentradas (como HEAs, RHEAs, CCAs e MPEAs). A segunda categoria reúne 22 descritores técnicos associados a solicitações hipersônicas e condições extremas, incluindo propulsão \emph{scramjet}, configurações \emph{waverider}, ablação, choque térmico e ensaios em túneis de plasma. O algoritmo computacional gerou o produto cartesiano de todos os pares possíveis unidos pelo operador \texttt{AND}, articulando-os sequencialmente via \texttt{OR} sobre os campos de Título (\texttt{TI}), Resumo (\texttt{AB}) e Palavras-chave do Autor (\texttt{AK}). Essa parametrização filtrou com precisão cirúrgica um corpus consolidado de 192 artigos científicos altamente especializados, cujas listas completas de palavras estão documentadas no Apêndice~\ref{apendice:query_wos}.

\subsection{Higienização e Padronização dos Dados}\label{sec:elegibilidade_processamento}

Bancos de dados bibliográficos brutos apresentam inconsistências terminológicas e sobreposições autorais que exigem tratamento prévio rigoroso. Inicialmente, construiu-se um vocabulário controlado (\emph{thesaurus}) para unificar as múltiplas designações de materiais sob uma taxonomia estável, impedindo que variações acrônimas dispersassem núcleos de pesquisa semanticamente idênticos. 

Adicionalmente, enfrentou-se o desafio crítico da homonímia autoral decorrente da indexação abreviada de pesquisadores pelo padrão sobrenome-inicial. Nomes recorrentes na literatura asiática foram desambiguados manualmente através do cruzamento de identificadores digitais persistentes (ORCID e Scopus Author ID), análise do histórico de afiliações institucionais e correlação de redes de coautoria habitual. Definiram-se como critérios de elegibilidade a inclusão exclusiva de artigos originais e revisões por pares focados em comportamento mecânico, microestrutural ou oxidativo em temperaturas superiores a $1200~^\circ\mathrm{C}$ e condições de tocha ou plasma. Ligas diluídas convencionais ou estudos restritos a temperaturas moderadas foram descartados. A base depurada foi então processada de forma integrada no pacote \emph{Bibliometrix} (ambiente R) para extração de indicadores cronológicos e no software \emph{VOSviewer} para mapeamento e clusterização de redes colaborativas.
